\chapter[Introdução]{Introdução}
\label{cap:introducao}

\section{Contextualização}

Robôs móveis autônomos deixaram de ser exclusivos de laboratórios e hoje
aparecem em armazéns, hospitais, inspeção industrial e residências. Para
se deslocar sozinho em um ambiente, um robô precisa resolver três
problemas ligados entre si \cite{siegwart2011}. Ele precisa perceber o
ambiente e construir um mapa dele. Precisa saber onde está nesse mapa. E
precisa planejar e executar o movimento até o destino. Os dois primeiros
problemas são tratados em conjunto pelas técnicas de localização e
mapeamento simultâneos (SLAM). O terceiro depende do planejamento de
caminho e, na base de tudo, do controle dos motores, que transforma os
comandos de velocidade em movimento real das rodas.

O ROS~2 (\textit{Robot Operating System 2}) se tornou a base comum para
integrar essas partes \cite{macenski2022ros2}. Ele oferece drivers para
sensores comerciais, algoritmos prontos de SLAM e navegação e um modelo
de comunicação que permite distribuir o processamento entre vários
computadores. Com isso, um robô funcional pode ser montado a partir de
componentes acessíveis, como computadores de placa única, sensores LiDAR
de baixo custo e motores com encoder.

O controle de movimento continua sendo uma área com muitas alternativas.
O controlador PID é o mais usado na prática, pela simplicidade e pela
robustez \cite{astrom2006}. O controle preditivo baseado em modelo (MPC)
trata explicitamente as restrições do sistema, ao custo de mais
processamento \cite{camacho2007}. Já os algoritmos bioinspirados, como a
otimização por enxame de partículas (PSO), permitem sintonizar
controladores sem um modelo exato da planta \cite{kennedy1995pso}. Cada
abordagem tem vantagens e limitações que dependem do sistema controlado e
do hardware em que ela roda.

\section{Justificativa}

Grande parte das comparações entre algoritmos de controle é feita em
simulação. A simulação é útil, mas não reproduz bem efeitos presentes em
um robô real: zona morta e saturação dos motores, atrito, escorregamento
das rodas, ruído nos sensores, atrasos de comunicação e o limite de
processamento do computador embarcado. Esses efeitos podem mudar qual
algoritmo tem o melhor desempenho na prática.

Comparar algoritmos em um robô real, por outro lado, exige uma plataforma
em que o controle seja acessível em todos os níveis. Muitas plataformas
comerciais usam controladoras de motor que fazem o controle de velocidade
dentro da própria placa, o que impede trocar o algoritmo desse nível.
Além disso, para que a comparação seja justa, todos os algoritmos
precisam ser avaliados no mesmo hardware, com os mesmos ensaios e as
mesmas métricas. Para isso, o software deve permitir trocar o controlador
sem mexer no restante do sistema.

Este trabalho propõe o ARENA, uma plataforma robótica móvel construída com
componentes de custo acessível. Os drivers de motor não fazem controle
próprio, então o algoritmo de controle fica inteiramente no software. Do
ponto de vista da Engenharia de Software, o projeto exige uma arquitetura
modular, com interfaces bem definidas entre os drivers, a odometria, a
navegação e os controladores. Isso permite incluir novos algoritmos sem
alterar o resto do sistema. A plataforma também pode servir de base para
outros trabalhos e para o ensino de robótica e controle.

\section{Questões de pesquisa}

Este trabalho busca responder às seguintes questões:

\begin{enumerate}
	\item É possível construir, com componentes de baixo custo e ROS~2, um
	robô 4WD capaz de mapear um ambiente e navegar de forma autônoma até
	um ponto escolhido no mapa?
	\item Como estruturar o software do robô para que diferentes
	algoritmos de controle possam ser trocados e comparados nas mesmas
	condições?
	\item Como o PID clássico, o MPC e os algoritmos bioinspirados se
	comparam, em um robô real, quanto ao erro de seguimento, ao esforço de
	controle e ao custo computacional?
	\item Quais desses algoritmos podem ser executados no computador
	embarcado e quais precisam ser executados na estação de base?
\end{enumerate}

\section{Objetivos}

\subsection{Objetivo geral}

Desenvolver o ARENA, uma plataforma robótica móvel de tração nas quatro
rodas baseada em ROS~2, capaz de mapear um ambiente e navegar de forma
autônoma, e usá-la para comparar experimentalmente algoritmos de controle
(PID clássico, controle preditivo baseado em modelo e algoritmos
bioinspirados) sob as mesmas condições de ensaio.

\subsection{Objetivos específicos}

\begin{enumerate}
	\item Montar e integrar o hardware do robô: computador embarcado,
	motores com encoder, drivers de motor, IMU, LiDAR, câmera de
	profundidade e sistema de alimentação;
	\item Caracterizar experimentalmente os motores e os encoders e obter
	um modelo da dinâmica de cada lado do robô;
	\item Desenvolver o nó da base em ROS~2, com controle de velocidade das
	rodas por algoritmo intercambiável, cálculo da odometria e parada
	automática em caso de falha de comunicação;
	\item Integrar os sensores ao ROS~2 e combinar a odometria das rodas com
	a IMU por um filtro de Kalman estendido;
	\item Construir o mapa do ambiente com o \texttt{slam\_toolbox} e
	realizar a navegação autônoma até um ponto do mapa com o Nav2;
	\item Implementar os controladores PID clássico, MPC e bioinspirados
	nos níveis de velocidade das rodas e de seguimento de trajetória;
	\item Definir um conjunto padronizado de ensaios e métricas e comparar
	os algoritmos quanto ao desempenho e ao custo computacional;
	\item Documentar e disponibilizar o código-fonte e os procedimentos de
	montagem e de ensaio, para que a plataforma possa ser reproduzida e
	ampliada.
\end{enumerate}

\section{Organização do trabalho}

Este trabalho está organizado da seguinte forma. O
Capítulo~\ref{cap:hardware} descreve a arquitetura de hardware do robô:
os componentes escolhidos, o funcionamento de cada um e a forma como se
integram. O Capítulo~\ref{cap:software} descreve a arquitetura de
software: as ferramentas de validação dos componentes, a organização do
sistema em nós do ROS~2, o fluxo de dados entre o robô e a estação de
base e a estrutura usada para comparar os algoritmos de controle,
incluindo a metodologia dos ensaios.
% TODO: incluir aqui os capitulos que forem acrescentados (referencial
% teorico, resultados parciais, cronograma, consideracoes finais).
